
\documentclass{article}
\usepackage{colortbl}
\usepackage{makecell}
\usepackage{multirow}
\usepackage{supertabular}

\begin{document}

\newcounter{utterance}

\centering \large Interaction Transcript for game `bbh', experiment `dyck\_languages', episode 40 with reflexion{-}Qwen3.5{-}4B.
\vspace{24pt}

{ \footnotesize  \setcounter{utterance}{1}
\setlength{\tabcolsep}{0pt}
\begin{supertabular}{c@{$\;$}|p{.15\linewidth}@{}p{.15\linewidth}p{.15\linewidth}p{.15\linewidth}p{.15\linewidth}p{.15\linewidth}}
    \# & \multicolumn{2}{c}{Player} && \multicolumn{2}{c}{Game Master} \\
    \hline

    \theutterance \stepcounter{utterance}  
    & & & \multicolumn{4}{p{0.6\linewidth}}{
        \cellcolor[rgb]{0.9,0.9,0.9}{
            \makecell[{{p{\linewidth}}}]{
                \texttt{\tiny{[P1$\langle$GM]}}
                \texttt{You have attempted this same task before and did not succeed. These are the plans you wrote after those attempts.} \\
\texttt{Trial \#1: New plan: I failed to correctly close the nested structure because I miscounted the open brackets and ignored the specific order of the remaining open symbols. Instead of guessing the closing sequence, I will first count the total number of open symbols in the input string, then iterate through them from left to right to determine the exact order in which they must be closed. I will generate the closing sequence by reversing the order of the unique open symbols found in the input, ensuring that every opening bracket has its corresponding closing bracket immediately following it.} \\
\\ 
\\ 
\texttt{Correctly close a Dyck{-}n word.} \\
\\ 
\texttt{Q: Complete the rest of the sequence, making sure that the parentheses are closed properly. Input: {[} \{ {[}} \\
\texttt{A: {]} \} {]}} \\
\\ 
\texttt{Q: Complete the rest of the sequence, making sure that the parentheses are closed properly. Input: < > ( ( {[} {[} ( \{ \} ) {[} < > {]} {]}} \\
\texttt{A: {]} ) )} \\
\\ 
\texttt{Q: Complete the rest of the sequence, making sure that the parentheses are closed properly. Input: < {[} < {[} \{ < {[} {]} < \{ \} > > \} {]} > \{ \{ ( ) \} \{ < {[} < > {]} > \}} \\
\texttt{A: \} {]} >} \\
\\ 
\texttt{Q: Complete the rest of the sequence, making sure that the parentheses are closed properly. Input: \{ < \{ \} > \{ (} \\
\texttt{A:} \\
            }
        }
    }
     \\ \\

    \theutterance \stepcounter{utterance}  
    & \multicolumn{4}{p{0.6\linewidth}}{
        \cellcolor[rgb]{0.561,0.737,0.561}{
            \makecell[{{p{\linewidth}}}]{
                \texttt{\tiny{[P1$\rangle$GM]}}
                \texttt{\} > )} \\
            }
        }
    }
    & & \\ \\

    \theutterance \stepcounter{utterance}  
    & & & \multicolumn{2}{p{0.3\linewidth}}{
        \cellcolor[rgb]{0.9,0.9,0.9}{
            \makecell[{{p{\linewidth}}}]{
                \texttt{\tiny{[GM$|$GM]}}
                \texttt{Target: ) \} \}} \\
            }
        }
    }
    & & \\ \\

    \theutterance \stepcounter{utterance}  
    & & & \multicolumn{2}{p{0.3\linewidth}}{
        \cellcolor[rgb]{0.9,0.9,0.9}{
            \makecell[{{p{\linewidth}}}]{
                \texttt{\tiny{[GM$|$GM]}}
                \texttt{game\_result = LOSE} \\
            }
        }
    }
    & & \\ \\

    \theutterance \stepcounter{utterance}  
    & & \multicolumn{4}{p{0.6\linewidth}}{
        \cellcolor[rgb]{0.95,0.95,0.95}{
            \makecell[{{p{\linewidth}}}]{
                \texttt{\tiny{[P1$\leftrightarrow$P]}}
                \texttt{Post{-}episode Reflexion (not sent to GM):} \\
\texttt{New plan: I failed to correctly close the nested structure because I miscounted the open brackets and ignored the specific order of the remaining open symbols. Instead of guessing the closing sequence, I will first count the total number of open symbols in the input string, then iterate through them from left to right to determine the exact order in which they must be closed. I will generate the closing sequence by reversing the order of the unique open symbols found in the input, ensuring that every opening bracket has its corresponding closing bracket immediately following it.} \\
            }
        }
    }
    & \\ \\

\end{supertabular}
}

\end{document}
